\section{Closed piecewise-\texorpdfstring{$C^{1}$}{C¹} immersions and null-homology}\label{sec:curves}

This chapter introduces the curve types underlying the
Hungerbühler--Wasem residue theorem. The basic carrier is the
\emph{closed piecewise-$C^{1}$ immersion}: a closed curve in
$\mathbb{C}$ that is $C^{1}$ on each of finitely many closed
sub-intervals between partition points, with nowhere-vanishing
derivative. The second concept, \emph{null-homology in an open set}
$U$, is the topological condition required to apply the residue
theorem along such a curve.

\begin{definition}[Closed piecewise-$C^{1}$ immersion]
  \label{def:closed-pwc1-immersion}
  \lean{ClosedPwC1Immersion}
  \leanok
  Let $x \in E$ be a point of a real normed vector space $E$.
  A \emph{closed piecewise-$C^{1}$ immersion} based at $x$ is the
  following data:
  \begin{itemize}
    \item a continuous path $\gamma : [0,1] \to E$ with
      $\gamma(0) = \gamma(1) = x$;
    \item a finite \emph{closed partition}
      $0 = a_0 < a_1 < \cdots < a_n = 1$ of $[0,1]$;
    \item on every closed sub-interval $[a_k, a_{k+1}]$ between
      consecutive partition members, the (extended) path
      $\gamma : [a_k, a_{k+1}] \to E$ is of class $C^{1}$ and its
      one-sided within-derivative is nowhere zero:
      \[
        \dot{\gamma}|_{[a_k, a_{k+1}]}(t) \ne 0
        \quad \text{for every } t \in [a_k, a_{k+1}].
      \]
  \end{itemize}
  This is the paper-faithful definition from Hungerbühler--Wasem 2018,
  page 3. The non-vanishing-on-each-closed-piece convention forces a
  well-defined right/left limit of $\dot{\gamma}$ at every partition
  point and, in particular, two non-zero \emph{tangent vectors} at each
  corner of $\gamma$.
\end{definition}

\begin{definition}[Null-homologous closed immersion]
  \label{def:null-homologous}
  \lean{IsNullHomologous}
  \uses{def:closed-pwc1-immersion}
  \leanok
  Let $U \subseteq \mathbb{C}$ be an open set, $x \in \mathbb{C}$, and
  $\gamma$ a closed piecewise-$C^{1}$ immersion based at $x$ in the
  sense of \cref{def:closed-pwc1-immersion}. We say that $\gamma$ is
  \emph{null-homologous in $U$} when both
  \begin{itemize}
    \item the image of $\gamma$ lies in $U$, that is
      $\gamma(t) \in U$ for every $t \in [0,1]$, and
    \item for every point $z \in \mathbb{C} \setminus U$, the
      generalised winding number
      $\operatorname{gWN}(\gamma, z)$ (see
      \cref{def:generalized-winding-number}) is zero.
  \end{itemize}
  Equivalently, $\gamma$ has image in $U$ and winds zero times around
  every point of the complement.
\end{definition}
